\section{Introduction}\label{sec:intro}

Machine-learning-based network intrusion detection systems (NIDS) routinely report accuracies and F1 scores above 0.99 on public benchmarks, yet practitioners remain reluctant to deploy them~\cite{sommer2010outside,arp2022dos,apruzzese2023sok}. A large part of the gap is that benchmark numbers answer a question operators do not ask. A security operations centre (SOC) does not need a classifier that separates benign and attack flows in a random split of one capture. It needs a system that, on \emph{its own} network, with \emph{today's} attacks and with a \emph{fixed} number of analyst hours, puts the right flows in front of analysts and does not bury them in false alarms~\cite{alahmadi2022falsepositives}. Three properties matter for that goal. The detector must transfer to traffic that differs from its training data (cross-dataset or domain shift)~\cite{layeghy2023explainable,cantone2024cross}. It must not silently wave through attack families it has never seen (open-set or zero-day behaviour)~\cite{talpini2024enhancing,farrukh2023detecting}. And its confidence must mean something, so that thresholds and alert volumes can be planned (calibration)~\cite{guo2017calibration,pal2025calibrated}. Because attacks are rare, all three must be judged at realistic base rates, where even a small false-positive rate dominates the alert queue~\cite{axelsson2000base}.

Each of these properties has been studied, but almost always in isolation. Cross-dataset studies show that performance can drop sharply when a model trained on one benchmark is applied to another~\cite{layeghy2023explainable,cantone2024cross,xin2025crossdataset}; open-set studies hold out attack classes within one dataset~\cite{talpini2024enhancing,farrukh2023detecting,yu2025industrial}; calibration and uncertainty studies typically stay in-distribution or within a single dataset~\cite{pal2025calibrated}; and conformal methods for security have focused on malware concept drift~\cite{jordaney2017transcend,barbero2022transcending} or adversarial perturbations~\cite{li2026robust,yumlembam2024comprehensive}. A recent review of 29 empirical deep-learning NIDS studies published between 2024 and 2026 found that \emph{none} evaluated cross-dataset generalisation, rejection of unseen attack families and confidence reliability together~\cite{lew2026deployment}. The closest empirical work evaluates calibration and fixed-false-positive-rate alerting under cross-dataset transfer, but for binary detection only, without unseen families or selective and conformal mechanisms, and with a single seed~\cite{jin2026benchmark}.

This paper closes that gap with a joint, budget-aware evaluation and a simple triage mechanism that can be added to any detector. Our contributions are:
\begin{enumerate}
\item \textbf{A joint evaluation protocol.} We evaluate four standard model families on five public benchmarks under three regimes (in-domain, six directed cross-dataset pairs over a harmonised 77-feature CICFlowMeter space, and leave-one-family-out), with five seeds. We report discrimination, calibration, false-alarm inflation, base-rate precision and analyst-budget metrics side by side, so that the same model can be judged on all of them at once.
\item \textbf{Budgeted conformal triage.} We frame alert triage as testing the null hypothesis ``benign and typical of the training traffic'' and combine a detector-based and a novelty-based conformal p-value calibrated on benign flows only. The resulting queue has a distribution-free false-alarm guarantee under exchangeability, and its size can be fixed to the SOC budget from \emph{unlabeled} target traffic.
\item \textbf{Empirical findings.} We quantify how far in-domain calibration and false-alarm guarantees survive under each type of shift, how often unseen attack families are missed with high confidence, and when novelty evidence helps or hurts a budgeted review queue (Section~\ref{sec:results}).
\item \textbf{Reproducibility.} All code, dataset provenance (URLs, commit hashes and checksums) and raw result tables are released, and every number in this paper is produced by that code.
\end{enumerate}

The remainder of the paper is organised as follows. Section~\ref{sec:related} reviews related work. Section~\ref{sec:method} describes the data, protocols and triage method. Section~\ref{sec:results} presents the results, Section~\ref{sec:discussion} discusses implications and threats to validity, and Section~\ref{sec:conclusion} concludes.
